% GENERATED FILE. Edit canonical lessons, then run npm run generate:books.
% canonical-source-hash: fnv1a64:06a790018fa7e7a2
% canonical-lessons: PT-C477-choro, PT-C477-solucar, PT-C477-suspirar, PT-C477-suspiro, PT-C477-gargalhada

\chapter{Lágrimas e gargalhadas}
\label{ch:pt-lagrimas-e-gargalhadas}
\hlchaptermodality{477}

\begin{chapteropening}
\textbf{By the end of this chapter:} \emph{I can describe crying, sighing and laughing out loud.}

Complete the last lesson of chapter 477: I can describe crying, sighing and laughing out loud.
\end{chapteropening}

\section[o choro]{o choro — crying}

\label{lesson:PT-C477-choro}



\begin{quote}
Before the new one: say the Portuguese for to console, then the Portuguese for to get something off your chest.
\end{quote}

\subsection*{You'll want to know: o choro}
\textbf{o choro} — \textquotedblleft{}crying\textquotedblright{}.

\begin{quote}
Ouvia-se o choro de um bebé no comboio.
\end{quote}

\begin{quote}
\emph{You could hear a baby crying on the train.}
\end{quote}

The noun of \emph{chorar}.

\begin{cousinweb}
\emph{choro} comes from \emph{chorar} \textquotedblleft{}to cry\textquotedblright{}, which is inherited from Latin \emph{plōrāre} \textquotedblleft{}to wail\textquotedblright{} (Lewis \& Short).
\end{cousinweb}

\begin{grammarlens}[title={What you've built}]
One more word for telling, explaining and describing at B1.
\end{grammarlens}

\subsection*{Your turn}
\begin{itemize}
\raggedright
  \item \emph{Say it:} \emph{o choro}
  \item \emph{Say it:} \emph{o choro}, once more
  \item \emph{Say it:} a sentence of your own with \emph{o choro}
  \item \emph{Recall:} say the Portuguese for to console, then the Portuguese for to get something off your chest, then say \emph{o choro} again
\end{itemize}

\begin{tcolorbox}[breakable,colback=teal!4,colframe=teal!35!black,title={Before you move on}]
What does o choro mean? (\textquotedblleft{}Crying\textquotedblright{}.) Say it once more.
\end{tcolorbox}
\section[soluçar]{soluçar — to sob; to hiccup}

\label{lesson:PT-C477-solucar}



\begin{quote}
Before the new one: say the Portuguese for to get something off your chest, then the Portuguese for crying.
\end{quote}

\subsection*{You'll want to know: soluçar}
\textbf{soluçar} — \textquotedblleft{}to sob; to hiccup\textquotedblright{}.

\begin{quote}
A criança soluçava e não conseguia falar.
\end{quote}

\begin{quote}
\emph{The child was sobbing and could not speak.}
\end{quote}

\emph{O soluço} is a sob, and also the hiccups: \emph{tenho soluços}.

\begin{cousinweb}
\emph{soluçar} is built on \emph{soluço} \textquotedblleft{}a sob, a hiccup\textquotedblright{}, which is inherited, through Old Portuguese \emph{solluço}, from Vulgar Latin \emph{suggluttium} (EtymDB).
\end{cousinweb}

\begin{grammarlens}[title={What you've built}]
One more word for telling, explaining and describing at B1.
\end{grammarlens}

\subsection*{Your turn}
\begin{itemize}
\raggedright
  \item \emph{Say it:} \emph{soluçar}
  \item \emph{Say it:} \emph{soluçar}, once more
  \item \emph{Say it:} a sentence of your own with \emph{soluçar}
  \item \emph{Recall:} say the Portuguese for to get something off your chest, then the Portuguese for crying, then say \emph{soluçar} again
\end{itemize}

\begin{tcolorbox}[breakable,colback=teal!4,colframe=teal!35!black,title={Before you move on}]
What does soluçar mean? (\textquotedblleft{}To sob; to hiccup\textquotedblright{}.) Say it once more.
\end{tcolorbox}
\section[suspirar]{suspirar — to sigh}

\label{lesson:PT-C477-suspirar}



\begin{quote}
Before the new one: say the Portuguese for crying, then the Portuguese for to sob.
\end{quote}

\subsection*{You'll want to know: suspirar}
\textbf{suspirar} — \textquotedblleft{}to sigh\textquotedblright{}.

\begin{quote}
O meu pai suspira sempre que vê a conta da luz.
\end{quote}

\begin{quote}
\emph{My father sighs every time he sees the electricity bill.}
\end{quote}

\emph{Suspirar por} — to long for.

\begin{cousinweb}
\emph{suspirar} is inherited, through Old Portuguese \emph{suspirar} (EtymDB), from Latin \emph{suspīrāre} \textquotedblleft{}to draw a deep breath, to sigh\textquotedblright{} (Lewis \& Short), built on \emph{spīrāre} \textquotedblleft{}to breathe\textquotedblright{}.
\end{cousinweb}

\begin{grammarlens}[title={What you've built}]
One more word for telling, explaining and describing at B1.
\end{grammarlens}

\subsection*{Your turn}
\begin{itemize}
\raggedright
  \item \emph{Say it:} \emph{suspirar}
  \item \emph{Say it:} \emph{suspirar}, once more
  \item \emph{Say it:} a sentence of your own with \emph{suspirar}
  \item \emph{Recall:} say the Portuguese for crying, then the Portuguese for to sob, then say \emph{suspirar} again
\end{itemize}

\begin{tcolorbox}[breakable,colback=teal!4,colframe=teal!35!black,title={Before you move on}]
What does suspirar mean? (\textquotedblleft{}To sigh\textquotedblright{}.) Say it once more.
\end{tcolorbox}
\section[o suspiro]{o suspiro — sigh; meringue}

\label{lesson:PT-C477-suspiro}



\begin{quote}
Before the new one: say the Portuguese for to sob, then the Portuguese for to sigh.
\end{quote}

\subsection*{You'll want to know: o suspiro}
\textbf{o suspiro} — \textquotedblleft{}sigh; meringue\textquotedblright{}.

\begin{quote}
Ela deu um suspiro e fechou o livro.
\end{quote}

\begin{quote}
\emph{She gave a sigh and closed the book.}
\end{quote}

The noun of \emph{suspirar}. In a Portuguese pastry shop, \emph{um suspiro} is a meringue.

\begin{cousinweb}
\emph{suspiro} comes from \emph{suspirar} \textquotedblleft{}to sigh\textquotedblright{}, inherited from Latin \emph{suspīrāre} (Lewis \& Short), from \emph{spīrāre} \textquotedblleft{}to breathe\textquotedblright{}.
\end{cousinweb}

\begin{grammarlens}[title={What you've built}]
One more word for telling, explaining and describing at B1.
\end{grammarlens}

\subsection*{Your turn}
\begin{itemize}
\raggedright
  \item \emph{Say it:} \emph{o suspiro}
  \item \emph{Say it:} \emph{o suspiro}, once more
  \item \emph{Say it:} a sentence of your own with \emph{o suspiro}
  \item \emph{Recall:} say the Portuguese for to sob, then the Portuguese for to sigh, then say \emph{o suspiro} again
\end{itemize}

\begin{tcolorbox}[breakable,colback=teal!4,colframe=teal!35!black,title={Before you move on}]
What does o suspiro mean? (\textquotedblleft{}Sigh; meringue\textquotedblright{}.) Say it once more.
\end{tcolorbox}
\section[a gargalhada]{a gargalhada — burst of laughter, guffaw}

\label{lesson:PT-C477-gargalhada}



\begin{quote}
Before the new one: say the Portuguese for to sigh, then the Portuguese for sigh.
\end{quote}

\subsection*{You'll want to know: a gargalhada}
\textbf{a gargalhada} — \textquotedblleft{}burst of laughter, guffaw\textquotedblright{}.

\begin{quote}
O público deu uma gargalhada.
\end{quote}

\begin{quote}
\emph{The audience burst out laughing.}
\end{quote}

\emph{Dar uma gargalhada} — to burst out laughing.

\begin{cousinweb}
\emph{gargalhada} is usually taken to imitate the sound of loud laughter in the throat, like \emph{gargarejar} \textquotedblleft{}to gargle\textquotedblright{}; the exact origin is not settled.
\end{cousinweb}

\begin{grammarlens}[title={What you've built}]
One more word for telling, explaining and describing at B1.
\end{grammarlens}

\subsection*{Your turn}
\begin{itemize}
\raggedright
  \item \emph{Say it:} \emph{a gargalhada}
  \item \emph{Say it:} \emph{a gargalhada}, once more
  \item \emph{Say it:} a sentence of your own with \emph{a gargalhada}
  \item \emph{Recall:} say the Portuguese for to sigh, then the Portuguese for sigh, then say \emph{a gargalhada} again
\end{itemize}

\begin{tcolorbox}[breakable,colback=teal!4,colframe=teal!35!black,title={Before you move on}]
What does a gargalhada mean? (\textquotedblleft{}Burst of laughter, guffaw\textquotedblright{}.) Say it once more.
\end{tcolorbox}
